% !PRINTED 70
\weberpage{70}
Es gilt nun der Satz:

\gesperrt{3. Sind $\Omega'$ und $\Omega''$ zwei Theiler von $\Omega_m$, die zu den
Gruppen $\mathfrak{A}'$ und $\mathfrak{A}''$ gehören, so ist $\Omega''$ dann und
nur dann ein Theiler von $\Omega'$, wenn $\mathfrak{A}'$ ein Theiler
von $\mathfrak{A}''$ ist.}

Denn wenn $\Omega''$ ein Theiler von $\Omega'$ ist, so müssen alle Zahlen
von $\Omega''$ die Substitutionen der Gruppe $\mathfrak{A}'$, zu der $\Omega'$ gehört, gestatten, also muss $\mathfrak{A}'$ in $\mathfrak{A}''$ enthalten sein. Und wenn umgekehrt $\mathfrak{A}'$ in $\mathfrak{A}''$
enthalten ist, so gestatten alle Zahlen von $\Omega''$
die sämmtlichen Substitutionen von $\mathfrak{A}'$ und sind also rational
durch eine primitive Zahl des Körpers $\Omega'$ darstellbar; d. h. $\Omega''$
ist in $\Omega'$ enthalten (§. 144 des ersten Bandes).

Der Körper $\Omega_m$ enthält als Theiler alle Körper $\Omega_{m_1}$, bei
denen $m_1$ ein Theiler von $m$ ist; denn $\Omega_{m_1}$ besteht aus allen
rationalen Functionen von $r^{\frac{m}{m_1}}$. Nun ist dann und nur dann
\[
r^{a\frac{m}{m_1}} = r^{\frac{m}{m_1}},
\]
wenn
\weq{4}{a \equiv 1\,(\mathrm{mod}\ m_1).}

Die Zahlen $a$ aus $\mathfrak{N}$, die der Congruenz (4) genügen, bilden
also die Gruppe, zu der der Körper $\Omega_{m_1}$ gehört. Wir wollen
diese Gruppe mit $\mathfrak{A}_{m_1}$ bezeichnen und symbolisch
\weq{5}{\mathfrak{A}_{m_1} \equiv 1\,(\mathrm{mod}\ m_1).}

setzen. Da der Grad des Körpers $\Omega_{m_1}$ gleich $\varphi(m_1)$ ist, was
zugleich der Index des Theilers $\mathfrak{A}_{m_1}$ von $\mathfrak{N}$ ist, so ist der Grad
von $\mathfrak{A}_{m_1}$ gleich $\varphi(m):\varphi(m_1)$.

Ist $m = m_1 m_2$ und $m_1$ relativ prim zu $m_2$, so erhält man
alle Zahlen von $\mathfrak{A}_{m_1}$ aus $a = 1 + x m_1$, wenn man $x$ so bestimmt,
dass $a$ nach dem Modul $m_2$ mit allen relativen Primzahlen zu
$m_2$ congruent wird. Daraus ergiebt sich in Uebereinstimmung
mit der obigen allgemeinen Bestimmung der Grad von $\mathfrak{A}_{m_1}$ für
diesen Fall $=\varphi(m_2)$.

Sind nun aber $m_1$ und $m_2$ irgend zwei Theiler von $m$, und
$\mu$ der grösste gemeinschaftliche Theiler von $m_1$ und $m_2$, so wird
der grösste gemeinschaftliche Theiler der beiden Körper $\Omega_{m_1}, \Omega_{m_2}$
zu der Gruppe gehören, die man als kleinstes gemeinschaftliches
Vielfaches der beiden Gruppen
\weq{6}{\mathfrak{A}_{m_1} \equiv 1\,(\mathrm{mod}\ m_1),\quad \mathfrak{A}_{m_2} \equiv 1\,(\mathrm{mod}\ m_2)}
%CONT>>
